\documentclass[preprint,12pt]{elsarticle}

\usepackage[T1]{fontenc}
\usepackage[utf8]{inputenc}
\usepackage{amsmath,amssymb}
\usepackage{booktabs}
\usepackage{tabularx}
\usepackage{array}
\usepackage{microtype}
\usepackage[hidelinks]{hyperref}

\journal{Knowledge-Based Systems}

\begin{document}

\begin{frontmatter}

\title{When Does Relationship Memory Help? Base-Relative Evidence Valuation for Live-Streaming Recommendation}

\author{Anonymous Authors}

\begin{keyword}
live-streaming recommendation \sep sequential recommendation \sep relationship memory \sep evidence valuation \sep selective prediction
\end{keyword}

\end{frontmatter}

\section{Introduction}

Sequential recommender systems infer users' evolving preferences from ordered interaction histories, and self-attentive models such as SASRec and BERT4Rec established strong sequence encoders for next-item recommendation \cite{kang2018sasrec,sun2019bert4rec}. Yet stronger sequence modeling does not remove a more basic question: when does an additional source of historical evidence still contain information that is incremental relative to what the base recommender already represents? More evidence is not necessarily more useful. An auxiliary source can be redundant when the base already represents overlapping information, ineffective when no target-specific evidence exists, or complementary when it retains useful information outside the base's current representation. The relevant problem is therefore not whether more historical evidence can improve recommendation in general, but whether that evidence has positive marginal value relative to what the base already captures.

Recent recommendation research has developed increasingly sophisticated mechanisms for integrating heterogeneous evidence, including adaptive gating across collaborative, semantic, structural, and long-/short-term preference signals and information-theoretic fusion that explicitly seeks complementary multimodal cues while reducing redundancy \cite{xie2026mhgfcl,zhu2026mifusr}. These advances primarily address how available evidence should be represented, weighted, aligned, or fused. This leaves a different methodological question: does a particular auxiliary evidence source contribute positive marginal value beyond a strong base recommender for the current recommendation event? We refer to this perspective as \emph{Base-relative evidence valuation}. The unit of analysis is not the standalone quality of an auxiliary model, but the incremental utility of its evidence conditional on what the base already represents.

Live-streaming recommendation provides a useful setting in which to study this question because candidate availability changes over time and repeat consumption of specific channels is substantial \cite{rappaz2021liverec}. KuaiLive further provides a real-time interactive benchmark with live-room timing information and multiple forms of user--streamer interaction, enabling evaluation under dynamic live-streaming conditions \cite{qu2026kuailive}. We therefore use a relationship-memory specialist as a transparent source of persistent user--creator evidence. For a target creator, this evidence can occupy three Base-relative analysis states: \emph{represented}, when the target relationship appears in the base-visible sequence; \emph{recoverable but unrepresented}, when it exists only in older history; or \emph{unavailable}, when no previous target-specific relationship exists. These states describe how the specialist's evidence relates to the current base representation; they are used only for mechanism analysis and are not available to the online selector.

Our empirical results show why this distinction matters. On KuaiLive, the aggregate ordering between the relationship-memory specialist and a strong base recommender reverses across candidate regimes even when the evaluated events and evidence-state composition are unchanged, indicating that regime dependence cannot be reduced to state prevalence alone. On Twitch/LiveRec, frozen development analysis reveals strong conditional heterogeneity across the three evidence states, and the same sign structure reproduces on untouched test data. A controlled development-only context-capacity intervention provides stronger mechanistic support: for the same events that move from recoverable-but-unrepresented to represented as the base context expands, specialist-minus-base utility falls by approximately 0.42--0.45 NDCG@10 while the canonical relationship-memory specialist is held fixed. This pattern supports a Base-relative visibility interpretation without implying a universal causal law, because operating regime and finer recency structure also modulate realized utility.

This heterogeneity motivates selective specialist use. We treat the event-level ranking-utility difference between the relationship-memory specialist and the base recommender as the decision quantity to be predicted, rather than relying only on where the base is weak. The distinction is substantive: base difficulty identifies hard events, whereas specialist-relative utility identifies events that this particular evidence source can improve. A difficult event may still contain no target-specific evidence recoverable by the relationship-memory specialist. In matched-budget analysis, utility-based selection concentrates more strongly on recoverable-but-unrepresented events, whereas a generic difficulty control disproportionately selects unavailable events. Post-hoc deferral and modern Learning-to-Defer already establish how routing rules can be learned on top of fixed predictors under accuracy--cost trade-offs \cite{narasimhan2022posthoc,mao2025mastering}. Our contribution is therefore not a new general deferral or gating formulation, but the characterization and operational use of the structured conditional utility of a transparent relationship-evidence specialist in recommendation.

The contributions of this work are threefold:
\begin{itemize}
    \item \textbf{Base-relative evidence valuation.} We formulate the value of persistent relationship evidence as its event-level marginal ranking utility relative to a strong base recommender, rather than as standalone specialist performance or generic base difficulty.
    \item \textbf{Relationship-evidence mechanism.} We characterize represented, recoverable-but-unrepresented, and unavailable evidence states, combine frozen development-to-test replication with a controlled base-context intervention, and show that evidence state structures---but does not deterministically fix---specialist utility across recommendation regimes.
    \item \textbf{Conditional decision and validation.} We connect specialist-relative utility to selective invocation and evaluate the resulting decision principle across KuaiLive regimes, strong-base controls, independent training seeds, and a frozen second-platform test.
\end{itemize}

\section{Related Work}

\subsection{Sequential and Long-Horizon Recommendation}

Sequential recommendation models evolving user preference from ordered interaction histories. Self-attentive architectures such as SASRec and BERT4Rec established strong Transformer-style baselines by learning dependencies within historical sequences \cite{kang2018sasrec,sun2019bert4rec}. Subsequent research has expanded the range and structure of historical information available to sequential recommenders through richer context modeling, auxiliary signals, and more expressive sequence representations. This literature establishes that broader history can improve recommendation quality. The present study assumes a strong sequential base and asks a different question: whether a separately maintained source of persistent relationship evidence retains positive marginal value after accounting for what the base already represents.

\subsection{Auxiliary Evidence Integration, Reliability, and Harmful Information}

A closely related research line studies how heterogeneous auxiliary evidence should be incorporated into recommendation. Recent KBS work illustrates several increasingly refined strategies. Xie et al. combine collaborative, semantic, and structural views through multi-head gating and explicitly distinguish users' long-term interests from short-term preferences \cite{xie2026mhgfcl}. Zhu et al. formulate multimodal sequential fusion through mutual information, seeking complementary cross-modal information while reducing redundancy and modality noise \cite{zhu2026mifusr}. Zhang et al. further address irrelevant and imbalanced multimodal content through collaborative denoising and modality balancing \cite{zhang2026mcd4sr}. CollaPPA integrates collaborative knowledge with personalized sequential preference alignment in an LLM-based recommender \cite{yue2026collappa}. These studies primarily improve the joint representation learned from multiple available sources.

A second group of methods focuses on the quality or reliability of auxiliary evidence. Huo et al. model content-level and interaction-level reliability in multimodal recommendation \cite{huo2026reliability}. Related information-theoretic and knowledge-aware methods likewise distinguish informative signals from noisy or redundant ones. Reliability is therefore a natural conceptual neighbor of the present problem, but it is not equivalent to Base-relative marginal value: a signal may be reliable yet add little when the base already represents overlapping information.

Conditional evidence use has also been studied outside recommendation. Yerragunta et al. propose dynamic conditional gating for question answering over incomplete knowledge bases, controlling how structured KB evidence and textual evidence are aggregated \cite{yerragunta2026gating}. Negative-transfer research provides a complementary warning against assuming that additional information is uniformly beneficial; Peng et al. show that multimodal pretraining can harm downstream sequential recommendation and develop a re-learning strategy to mitigate that effect \cite{peng2025negative}. These studies motivate conditional treatment of auxiliary information without implying the same mechanism as the one studied here.

A broader decision-theoretic connection is the value of information (VOI), which quantifies the benefit that additional information can provide to a decision, often through expected utility improvement \cite{abbas2025voi}. Our setting differs from classical information acquisition: the auxiliary source is not an unobserved feature whose revelation must be purchased before acting, but a fixed relationship-memory specialist whose ranking contribution is evaluated relative to an already trained base recommender.

Collectively, these lines of work address the representation, reliability, conditional integration, or potential harm of auxiliary information. The present work takes event-level marginal ranking utility relative to a trained base recommender as the primary object of analysis; in this sense, Base-relative evidence valuation complements rather than replaces evidence-integration methods.

\subsection{Expert Routing, Selective Prediction, and Learning-to-Defer}

Selective prediction and Learning-to-Defer (L2D) provide the closest general decision-theoretic background. Mozannar and Sontag formulate L2D as jointly learning a predictor and a rejector that can defer decisions to an expert, with consistency results through cost-sensitive learning \cite{mozannar2020defer}. Narasimhan et al. study post-hoc deferral on top of a pretrained base model and optimize an accuracy--cost trade-off relative to an expert \cite{narasimhan2022posthoc}. Mao et al. develop stronger theoretical guarantees for single- and multiple-expert routing in both one-stage and two-stage settings \cite{mao2025mastering}, and recent work extends L2D to dynamically varying expert pools \cite{montreuil2026online}. Selective invocation of an expert, including on top of a fixed predictor, is therefore an established decision problem.

Recommendation also contains expert-style architectures. CeDRec, for example, uses adaptive in-context experts as components of a sequential representation \cite{sun2025cedrec}. Such systems adapt how representational experts contribute within the model, whereas the relationship-memory specialist studied here is maintained as a separate, transparent evidence source. The primary object of analysis is its specialist-minus-base ranking utility and the evidence conditions under which that utility changes. Accordingly, the generic difficulty control in our experiments serves a narrower mechanistic purpose: it tests whether base weakness alone is sufficient to identify cases where the relationship-memory specialist helps, rather than standing in for the broader L2D literature.

\subsection{Live-Streaming and Repeat-Aware Recommendation}

Live-streaming recommendation introduces domain-specific dynamics that are absent from static-item settings: rooms become available and unavailable over time, viewers repeatedly return to channels, and interactions between viewers and streamers evolve during sessions. Zhang et al. model the two-sided dynamics of viewers and anchors through bi-directional prediction \cite{zhang2021bidirectional}. Rappaz et al. identify dynamic availability and repeat consumption as defining characteristics of live-stream recommendation and propose LiveRec to model them \cite{rappaz2021liverec}. Later work develops dynamic room representation and interaction-aware preference modeling \cite{gao2023driver,chen2024preference}, while hypergraph and graph-based approaches capture richer viewer--channel interactions and repeated behaviors \cite{yu2024hypergraph,zhu2025repeated}.

These studies establish that dynamic availability, repeated consumption, and viewer--streamer interactions are central to live recommendation. Live streaming is therefore not introduced here as a new source of repeat behavior, but as an empirical setting in which target-specific relationship evidence can be represented in the base-visible sequence, remain recoverable only from older history, or be unavailable while the candidate environment itself changes. KuaiLive and Twitch/LiveRec provide complementary environments for studying this Base-relative evidence value \cite{qu2026kuailive,rappaz2021liverec}.

\subsection*{Positioning Summary}

Table~\ref{tab:positioning} summarizes the main distinction between Base-relative evidence valuation and neighboring research lines.

\begin{table}[htbp]
\centering
\caption{Positioning relative to neighboring research lines.}
\label{tab:positioning}
\small
\renewcommand{\arraystretch}{1.15}
\begin{tabularx}{\textwidth}{>{\raggedright\arraybackslash}p{0.25\textwidth} >{\raggedright\arraybackslash}p{0.32\textwidth} >{\raggedright\arraybackslash}X}
\toprule
Research line & Primary objective & Difference in focus \\
\midrule
Sequential / long-horizon recommendation & Encode broader or richer interaction history & Evaluate residual utility of a separate relationship-evidence specialist. \\
Auxiliary / multimodal fusion & Integrate heterogeneous information into a joint representation & Evaluate whether a source is incrementally useful relative to the base. \\
Reliability / denoising & Estimate which signals are trustworthy, informative, or noisy & Reliability does not imply positive Base-relative marginal utility. \\
Value of information & Quantify the expected decision benefit of acquiring information & Evaluate the marginal ranking contribution of a fixed specialist relative to a trained base. \\
L2D / expert routing & Assign instances to predictors or experts under loss or cost & Characterize the utility structure of one interpretable evidence specialist. \\
Live-stream recommendation & Model dynamic availability, interactions, and repeated behavior & Use the domain to study Base-relative evidence value rather than introduce repeat modeling itself. \\
\bottomrule
\end{tabularx}
\end{table}

\section{Problem Formulation}
\label{sec:problem}

We formalize Base-relative evidence valuation as an event-level ranking-utility problem. The formulation separates three objects that are conceptually distinct throughout the paper: the utility achieved by a fixed base recommender, the incremental utility contributed by a fixed relationship-memory specialist, and the decision of whether that specialist should be invoked for a particular recommendation event. This separation is important because the specialist's realized benefit is observable only after the relevant target is known, whereas the serving-time decision must be made from information available before that outcome is observed.

\subsection{Recommendation Setting and Event-Level Ranking Utility}
\label{sec:event-utility}

Consider a recommendation event
\begin{equation}
    x=(u,t,\mathcal{C}_t,\mathcal{H}_{u,t},y),
    \label{eq:event}
\end{equation}
where $u$ denotes the user, $t$ the recommendation time, $\mathcal{C}_t$ the set of candidates available at time $t$, $\mathcal{H}_{u,t}$ the interaction history available strictly before the event, and $y\in\mathcal{C}_t$ the observed relevant target. The explicit dependence of $\mathcal{C}_t$ on time accommodates the changing room availability that characterizes live-stream recommendation.

Let $A$ be any ranking model and let $r_A(x)$ denote the 1-indexed rank assigned to $y$ among $\mathcal{C}_t$. For a cutoff $K$, we define the event-level ranking utility as
\begin{equation}
    u_K(A,x)=
    \begin{cases}
        \displaystyle \frac{1}{\log_2\!\left(r_A(x)+1\right)}, & r_A(x)\le K,\\[6pt]
        0, & r_A(x)>K.
    \end{cases}
    \label{eq:event-ndcg}
\end{equation}
With a single relevant target per event, Eq.~\eqref{eq:event-ndcg} is exactly the event-level contribution to NDCG@$K$. The aggregate ranking utility of $A$ over a distribution of recommendation events is therefore
\begin{equation}
    U_K(A)=\mathbb{E}\!\left[u_K(A,X)\right],
    \label{eq:aggregate-utility}
\end{equation}
and its empirical counterpart is the mean of Eq.~\eqref{eq:event-ndcg} over the evaluation events. Other event-level ranking utilities, such as Hit@$K$, can be substituted without changing the decision formulation below.

\subsection{Base-Relative Specialist Utility}
\label{sec:relative-utility}

Let $B$ denote a fixed base recommender and $M$ a fixed relationship-memory specialist. For any given selective-decision experiment, both ranking functions are fixed before the invocation policy is learned. We define the specialist-minus-base utility of event $x$ as
\begin{equation}
    \Delta_m(x)=u_K(M,x)-u_K(B,x).
    \label{eq:delta-memory}
\end{equation}
Thus, $\Delta_m(x)>0$ indicates that invoking the specialist improves event-level ranking utility relative to the base, $\Delta_m(x)=0$ indicates no marginal ranking benefit, and $\Delta_m(x)<0$ indicates negative marginal value relative to the base. The relevant quantity is therefore not the standalone quality of $M$, but its improvement over the already available decision of $B$.

To make this connection explicit, let $a(x)\in\{0,1\}$ denote a selective action, where $a(x)=0$ uses the base ranking and $a(x)=1$ uses the specialist ranking. The event-level utility of the resulting selective system is
\begin{align}
    u_K^{\mathrm{sel}}(x;a)
    &=\left(1-a(x)\right)u_K(B,x)+a(x)u_K(M,x)\\
    &=u_K(B,x)+a(x)\Delta_m(x).
    \label{eq:selective-identity}
\end{align}
Consequently, the expected improvement over always using the base is
\begin{equation}
    \mathbb{E}\!\left[u_K^{\mathrm{sel}}(X;a)-u_K(B,X)\right]
    =\mathbb{E}\!\left[a(X)\Delta_m(X)\right].
    \label{eq:expected-selective-gain}
\end{equation}
Equation~\eqref{eq:expected-selective-gain} gives the decision rationale for learning specialist-relative utility: under a selective invocation policy, only the utility differences on invoked events contribute to improvement over the base.

The realized $\Delta_m(x)$ is an offline supervisory quantity. Computing it requires the observed target $y$ and the rankings produced by both $B$ and $M$; it is therefore unavailable when the serving-time decision must be made. The practical problem is to predict its conditional expectation using information observable before the target outcome is known.

\subsection{Conditional Utility Estimation and Observable Information}
\label{sec:conditional-utility}

Let $Z$ denote the information available at the time of the invocation decision. In the general formulation, $Z$ may contain pre-outcome user and history summaries, candidate-set descriptors, and confidence information produced by the base recommender. The conditional specialist utility is
\begin{equation}
    \eta(z)=\mathbb{E}\!\left[\Delta_m(X)\mid Z=z\right],
    \label{eq:eta}
\end{equation}
and a learned estimator $\hat{\eta}(z)=g_{\theta}(z)$ is used to rank or classify events according to their expected specialist-minus-base value.

The information boundary is central to the formulation. The observed target $y$, realized ranks, $\Delta_m(x)$, and the Base-relative evidence states introduced later for mechanism analysis are offline quantities and are not included in $Z$. In particular, the represented, recoverable-but-unrepresented, and unavailable states are diagnostic concepts used to explain specialist behavior rather than oracle inputs to the selector.

This also clarifies why generic base difficulty and specialist utility are not equivalent decision targets. A difficulty score concerns the performance or uncertainty of $B$ alone, whereas $\eta(z)$ depends on the relative performance of both $M$ and $B$. A hard event for the base may remain hard for the specialist, so base weakness does not imply positive specialist-minus-base utility.

\subsection{Selective Specialist Use under Cost and Budget Constraints}
\label{sec:decision-rule}

The deployment decision may additionally account for the incremental cost of invoking the specialist. Let $c(X)\ge 0$ denote this cost and let $a(Z)\in\{0,1\}$ be a policy based only on serving-time information. For a trade-off parameter $\lambda\ge 0$, consider the expected net utility
\begin{equation}
    \mathcal{J}(a)
    =\mathbb{E}\!\left[a(Z)\Delta_m(X)-\lambda a(Z)c(X)\right].
    \label{eq:cost-objective}
\end{equation}
Defining
\begin{equation}
    \kappa(z)=\mathbb{E}\!\left[c(X)\mid Z=z\right],
\end{equation}
the pointwise expected-utility decision is
\begin{equation}
    a^{*}(z)
    =\mathbf{1}\!\left[\eta(z)>\lambda\kappa(z)\right].
    \label{eq:cost-rule}
\end{equation}
When the incremental specialist cost is approximately constant, Eq.~\eqref{eq:cost-rule} reduces to thresholding conditional utility at a fixed cost-adjusted level. We use this formulation as a decision-theoretic interpretation of selective specialist use rather than as a claim of a new general deferral rule.

For controlled evaluation, we additionally consider an exact invocation budget. Suppose an evaluation set contains $n$ events and the specialist may be invoked on exactly $B_q$ of them. Given predicted utilities $\hat{\eta}(z_i)$, the budget-constrained decision problem is
\begin{equation}
    \max_{a_1,\ldots,a_n\in\{0,1\}}
    \sum_{i=1}^{n}a_i\hat{\eta}(z_i)
    \quad
    \text{s.t.}\quad
    \sum_{i=1}^{n}a_i=B_q.
    \label{eq:budget-objective}
\end{equation}
With equal per-invocation cost, the solution is to select the $B_q$ events with the largest predicted utilities,
\begin{equation}
    a_{B_q}
    =\operatorname{Top}_{B_q}\!\left\{\hat{\eta}(z_i)\right\}_{i=1}^{n}.
    \label{eq:top-budget}
\end{equation}
This exact-budget form provides the bridge between the general cost-sensitive decision in Eq.~\eqref{eq:cost-rule} and the matched-budget comparisons used in the experiments.

\subsection{Regime-Level Utility Decomposition}
\label{sec:regime-decomposition}

Finally, we formalize how aggregate specialist value can change across operating recommendation regimes. Let $R$ denote a regime, such as a candidate-construction or temporal-availability regime, and let $S$ denote a Base-relative relationship-evidence state. The state variable is used for analysis and is instantiated operationally in Section~4; it need not be an online selector feature. By the law of total expectation,
\begin{equation}
    \mathbb{E}\!\left[\Delta_m(X)\mid R=r\right]
    =\sum_{s}P(S=s\mid R=r)
    \,\mathbb{E}\!\left[\Delta_m(X)\mid S=s,R=r\right].
    \label{eq:regime-decomposition}
\end{equation}
Equation~\eqref{eq:regime-decomposition} separates two distinct sources of aggregate regime dependence: changes in the prevalence of evidence states, $P(S=s\mid R=r)$, and changes in state-specific specialist utility, $\mathbb{E}[\Delta_m\mid S=s,R=r]$. We use this identity as an analytical decomposition rather than as a new theoretical result. In experiments conducted within a fixed operating regime, $R$ is implicit in the event distribution and is not assumed to be provided as an additional serving-time feature.

Together, Eqs.~\eqref{eq:delta-memory}--\eqref{eq:regime-decomposition} define the central objects used throughout the remainder of the paper: event-level Base-relative evidence value, its serving-time conditional expectation, the corresponding selective decision, and the decomposition of aggregate specialist value across recommendation regimes. Section~4 instantiates these objects with the two sequential base recommenders, the transparent relationship-memory specialist, observable selector features, and the diagnostic evidence states used in our mechanism analysis.

\bibliographystyle{elsarticle-num}
\bibliography{kbs_references}

\end{document}
